% ID: SEL-D0017-BB-E03
% ESTADO: cotejada. Segunda pasada visual completa y borrador compilado revisados.
% Fuente: Selectividad/matematicas_aplicadas/matematicas_aplicadas_P_E.pdf
% SHA-256: bc88bfeb0f184eb0c806c4efe04e84c24f3e5fcef9caaa12e94fd48d9f67a389
% Página: 1; bloque B; ejercicio original 3.

Se considera la función
\[
f(x)=\begin{cases}
x^2+ax+2 & \text{si }x\leq 0\\
\dfrac{x+b}{x-1} & \text{si }x>0
\end{cases}
\]
\begin{enumerate}[label=\alph*)]
  \item \textbf{(1.2 puntos)} Halle $a$ y $b$ para que $f$ sea continua
  y derivable en $x=0$.
  \item \textbf{(0.7 puntos)} Para $a=1$ y $b=-2$, halle la ecuación de la
  recta tangente a la gráfica de $f$ en el punto de abscisa $x=0$.
  \item \textbf{(0.6 puntos)} Para $a=1$ y $b=1$, halle, si existen, las
  ecuaciones de las asíntotas de $f$.
\end{enumerate}
